Energy efficiency has emerged as a vital attribute of software quality, with significant implications for both environmental sustainability and operational costs. However, existing profiling tools operate only at runtime and coarse granularity, typically capturing energy at the process or method level. Such tools fail to expose how small code blocks, such as functions, loops, and conditionals, contribute to energy consumption, preventing developers from reasoning about and comparing the energy efficiency of programming constructs during design-time. 

To address this gap, we propose EnCoDe, a methodology for fine-grained, design-time energy estimation, with the following key contributions: (1) PowerLens, a novel measurement methodology that achieves reliable sub-millisecond energy readings for small code blocks; (2) Extensive empirical study on code blocks extracted from over 18,000 Python programs, uncovering linear and non-linear relationships between energy consumption and static code features such as structural, complexity, density, and contextual characteristics, resulting in a first-of-its-kind fine-grained dataset; and (3) Predictive modeling, in which machine learning models are trained on these features to accurately estimate and classify block-level energy consumption at design-time. Our results demonstrate stable, reproducible block-level estimations, with regressors achieving $R^2 = 0.75$ and classifiers achieving $80.6\%$ accuracy in identifying energy hotspots, enabling developers to localize and address inefficient code regions early in the development process without execution.

%Our evaluation demonstrates that EnCoDe provides stable and reproducible energy estimations at the block level. By effectively identifying energy hotspots during design-time, the methodology enables developers to localize and mitigate inefficient code regions early in the development lifecycle without the need for program execution.





%The increasing reliance of software across all domains of life has amplified its contribution on global energy consumption. Despite growing awareness, energy-aware software development remains uncommon, as developers lack practical mechanisms to assess or optimize the energy footprint of their code during design time. Existing profiling tools measure energy only at runtime and at coarse granularity, which prevents the identification of specific energy hotspots and makes optimization harder. To address these challenges, we propose \textbf{GreenGauge}, a fine-grained, design-time framework for estimating the energy consumption of software.
%At the core of GreenGauge lies \textbf{PowerLens}, a measurement approach that overcomes the sampling limitations of hardware interfaces such as RAPL to capture sub-millisecond energy readings for individual code blocks, including functions, loops, and conditionals. These measurements serve as ground truth for predictive modeling based on static code metrics and structural features derived from Abstract Syntax Trees (ASTs) that capture distinct aspects of program behavior.
%The correlation analysis on code properties and energy usage reveals the important features which determine energy use. To provide insights, GreenGauge employs regression models to predict precise energy estimates and classification models to provide intuitive categorical feedback (low, medium, high) to guide developers during design and refactoring.
%By integrating these capabilities directly into the development environment, GreenGauge enables developers to make energy-conscious design decisions early in the software lifecycle, elevating energy efficiency to a first-class software quality attribute.


